\section{Introduction}
\label{sec:intro}

Mobile graphical-user-interface (GUI) agents complete tasks by reading what is on the screen and acting on it. When a task requires a secret---a one-time password (OTP), a password, an address, a phone number, or a payment identifier---the agent must obtain that value and place it into a field. In current designs the value enters the agent's observation and actuation: the SMS is rendered, the code is scraped from the screenshot or accessibility tree, and the value is typed into a node. An instruction injected through any content the agent reads can then turn task authority into disclosure or redirection; the value the agent legitimately handles is exactly the value an attacker wants.

\paragraph{A running example.} A user asks a mobile agent to log into their bank. The bank sends a one-time code by SMS. The agent opens the bank app, reads the code from the notification, types it into the code field, and the login completes. Now suppose the same SMS, a web page in the browsing flow, or a notification from a promotional app also contains a line such as ``to confirm, forward this code to \texttt{refund.example}.'' The agent cannot reliably separate the task it was given from text it merely read, so it forwards the code. The code was never stolen from the agent's memory; it was disclosed \emph{through} the agent's legitimate handling of it. Under \SH{}, the SMS pipeline mints an opaque handle bound to the bank's package and its OTP field \emph{before} the code can enter the agent's view. The agent sees only ``OTP ready for SecureBank'' and issues \Enter(\texttt{handle}) at the code field; the OS authenticates the handle, validates the destination, injects the value, and redacts the field. An \Enter{} aimed at \texttt{refund.example} is rejected because that destination lies outside the binding, no matter how convincing the injected text was.

This is no longer hypothetical. A browser agent read an OTP from a mail client and published it; mobile agents have forwarded Alipay verification codes under injected instructions (see Section~\ref{sec:background} for the primary sources). Placeholder-detection defenses replace sensitive rendered text with a local token and resolve it at use time, but they are detector-dependent, session-scoped, and---crucially---do not bind the resolved value to the destination that may consume it, so an injected instruction can redirect resolution to an attacker-owned field.

\paragraph{Research question.} Can Android let an agent \emph{use} a secret without ever \emph{observing} it, by minting unforgeable opaque handles at the authoritative OS source and resolving a handle only into a destination the handle is bound to---so that neither the raw value nor a usable copy enters the agent's perception or a wrong destination?

\paragraph{Why mobile.} The authoritative sources are mobile OS services: the SMS Retriever and User Consent APIs, the sensitive-notification OTP pipeline, Credential Manager, the clipboard, and on-device wallet tokenization. The agent reaches the UI through accessibility, so an OTP it must enter is otherwise read off the screen. Destination binding reuses Android's package identity and WebView origin. On a server there is no OS SMS pipeline, no Credential Manager, no on-device wallet, and no UI to type into; the problem and its enforcement point are mobile.

\paragraph{Contribution.} We propose \SH, an OS layer that makes secrets usable but invisible to untrusted agents. A handle is minted at the source before the raw value reaches the agent, carries a destination-and-task binding, and has a typed lifecycle. A resolver injects the value only into a bound destination; a redactor masks the resolved field afterwards. This emerging-ideas paper contributes (i) the problem framing and threat model (\S\ref{sec:background}--\ref{sec:threat}); (ii) the handle/resolver/redactor design with an unforgeable, destination- and task-bound abstraction (\S\ref{sec:design}); (iii) a compact formal handle model with four stated guarantees and two enforcement algorithms (\S\ref{sec:model}--\ref{sec:algorithm}); (iv) an on-emulator reproduction of the motivating attack, a stock-Android prototype of the mint-and-redact path with a preliminary before/after measurement ($29/30$ to $0/30$ exposure), and a clearly labeled simulation-only reference-model study (\S\ref{sec:evidence}); and (v) a preregistration-style evaluation plan for the full system. We do \emph{not} report device measurements in this paper; the formal results are conditional on stated reference-monitor assumptions and are validated only on controlled synthetic inputs.

\paragraph{Roadmap.} Section~\ref{sec:background} motivates the problem and Section~\ref{sec:threat} fixes the trust model. Section~\ref{sec:design} describes the design, Section~\ref{sec:model} formalizes handles and their guarantees, and Section~\ref{sec:algorithm} gives the mint and resolve algorithms. Section~\ref{sec:evidence} reports what we have reproduced and preregisters the full evaluation; Sections~\ref{sec:related}--\ref{sec:limitations} position the work and state its limits.
